\documentclass[reqno]{amsart}
\usepackage[T1]{fontenc}
\usepackage{amsmath,amssymb,amsthm}
\usepackage{xr-hyper}
\usepackage[hidelinks]{hyperref}
\externaldocument[sm-]{20261125T000000--slope-method__workbook}
\externaldocument[tt-]{20261125T004100--tate-twists-and-sha__workbook}
\externaldocument[oc-]{20261126T020430--the-minus-part-by-gauss-sums__workbook}
\externaldocument[cg-]{20261125T014100--congruence-modules__workbook}

% Chapter number; \ref{chap} gives it.
\newcounter{chap}
% Chapter numbers, generated from 20261125T012000--grh-workbook__guide.org
% by scripts/make-chapters.py.  Do not edit: change the order of the
% index and rerun the script (./preview-site.sh and CI do this).
%
% Usage in a sheet, after \newcounter{chap}:
%   % Chapter numbers, generated from 20261125T012000--grh-workbook__guide.org
% by scripts/make-chapters.py.  Do not edit: change the order of the
% index and rerun the script (./preview-site.sh and CI do this).
%
% Usage in a sheet, after \newcounter{chap}:
%   % Chapter numbers, generated from 20261125T012000--grh-workbook__guide.org
% by scripts/make-chapters.py.  Do not edit: change the order of the
% index and rerun the script (./preview-site.sh and CI do this).
%
% Usage in a sheet, after \newcounter{chap}:
%   \input{chapters}\chapnumber{<sheet ID>}
\makeatletter
\providecommand\chapters@set[2]{\@namedef{chapters@#1}{#2}}
\chapters@set{20261125T001000}{1}% zeros-to-ramp-sums
\chapters@set{20261125T001100}{2}% inert-primes-and-estimate-b
\chapters@set{20261125T005000}{3}% galois-fourier-and-averaging
\chapters@set{20261125T005500}{4}% columns-from-a-sphere
\chapters@set{20261125T000000}{5}% slope-method
\chapters@set{20261126T030100}{6}% a-twist-for-each-prime
\chapters@set{20261125T003000}{7}% the-chain-and-the-square-root
\chapters@set{20261125T010000}{8}% rank-lemma-and-zero-estimate
\chapters@set{20261125T010500}{9}% poincare-series-and-near-orthogonality
\chapters@set{20261124T223000}{10}% elliptic-curves-for-the-proof
\chapters@set{20261125T009500}{11}% fock-model-and-theta-intertwiners
\chapters@set{20261124T223100}{12}% complex-multiplication-and-inert-reduction
\chapters@set{20261125T008000}{13}% one-cm-determinant
\chapters@set{20261124T214500}{14}% supersingular-theta-divisibility
\chapters@set{20261124T214600}{15}% the-isotropic-plane
\chapters@set{20261125T013700}{16}% pairings-on-a-formal-group-of-height-two
\chapters@set{20261125T013500}{17}% hecke-nondegeneracy-on-the-tree
\chapters@set{20261126T030200}{18}% the-bound-is-attained
\chapters@set{20261125T009000}{19}% kronecker-limit-and-chowla-selberg
\chapters@set{20261125T002000}{20}% definite-theta-lifts
\chapters@set{20261125T015500}{21}% siegel-weil-in-the-smallest-case
\chapters@set{20261125T002100}{22}% definite-theta-lifts-and-norms
\chapters@set{20261125T016500}{23}% hecke-operators-and-theta-lifts
\chapters@set{20261125T016700}{24}% divisibility-along-the-chain
\chapters@set{20261125T004000}{25}% denominators-and-sha
\chapters@set{20261125T004100}{26}% tate-twists-and-sha
\chapters@set{20261126T020430}{27}% the-minus-part-by-gauss-sums
\chapters@set{20261125T014100}{28}% congruence-modules
\chapters@set{20261125T014500}{29}% eisenstein-congruences-and-selmer-classes
\chapters@set{20261125T014700}{30}% pseudocharacters-and-the-length-bound
\chapters@set{20261125T011000}{31}% galois-representations-black-box
\chapters@set{20261126T030300}{32}% the-data
\chapters@set{20261126T020210}{33}% bessel-poincare-series-and-their-gram-operators
\chapters@set{20261126T030400}{34}% the-seesaw-identity
\chapters@set{20261126T030500}{35}% divisibility-of-the-composite-lift
\chapters@set{20261126T030600}{36}% theorem-a-short-forms
\chapters@set{20261126T030700}{37}% the-determinant-of-a-chain
\chapters@set{20261126T030800}{38}% the-cost-of-one-lift
\chapters@set{20261126T030900}{39}% denominators-of-projections-to-lifts
\chapters@set{20261126T031000}{40}% twists-with-few-congruences
\chapters@set{20261126T031100}{41}% theorem-c-and-the-main-inequality
\providecommand\chapnumber[1]{%
  \@ifundefined{chapters@#1}%
    {\PackageError{chapters}{Sheet #1 is not listed in the index}%
      {List it in 20261125T012000--grh-workbook__guide.org\MessageBreak
       and run scripts/make-chapters.py.}}%
    {\edef\thechap{\@nameuse{chapters@#1}}}}
\makeatother
\chapnumber{<sheet ID>}
\makeatletter
\providecommand\chapters@set[2]{\@namedef{chapters@#1}{#2}}
\chapters@set{20261125T001000}{1}% zeros-to-ramp-sums
\chapters@set{20261125T001100}{2}% inert-primes-and-estimate-b
\chapters@set{20261125T005000}{3}% galois-fourier-and-averaging
\chapters@set{20261125T005500}{4}% columns-from-a-sphere
\chapters@set{20261125T000000}{5}% slope-method
\chapters@set{20261126T030100}{6}% a-twist-for-each-prime
\chapters@set{20261125T003000}{7}% the-chain-and-the-square-root
\chapters@set{20261125T010000}{8}% rank-lemma-and-zero-estimate
\chapters@set{20261125T010500}{9}% poincare-series-and-near-orthogonality
\chapters@set{20261124T223000}{10}% elliptic-curves-for-the-proof
\chapters@set{20261125T009500}{11}% fock-model-and-theta-intertwiners
\chapters@set{20261124T223100}{12}% complex-multiplication-and-inert-reduction
\chapters@set{20261125T008000}{13}% one-cm-determinant
\chapters@set{20261124T214500}{14}% supersingular-theta-divisibility
\chapters@set{20261124T214600}{15}% the-isotropic-plane
\chapters@set{20261125T013700}{16}% pairings-on-a-formal-group-of-height-two
\chapters@set{20261125T013500}{17}% hecke-nondegeneracy-on-the-tree
\chapters@set{20261126T030200}{18}% the-bound-is-attained
\chapters@set{20261125T009000}{19}% kronecker-limit-and-chowla-selberg
\chapters@set{20261125T002000}{20}% definite-theta-lifts
\chapters@set{20261125T015500}{21}% siegel-weil-in-the-smallest-case
\chapters@set{20261125T002100}{22}% definite-theta-lifts-and-norms
\chapters@set{20261125T016500}{23}% hecke-operators-and-theta-lifts
\chapters@set{20261125T016700}{24}% divisibility-along-the-chain
\chapters@set{20261125T004000}{25}% denominators-and-sha
\chapters@set{20261125T004100}{26}% tate-twists-and-sha
\chapters@set{20261126T020430}{27}% the-minus-part-by-gauss-sums
\chapters@set{20261125T014100}{28}% congruence-modules
\chapters@set{20261125T014500}{29}% eisenstein-congruences-and-selmer-classes
\chapters@set{20261125T014700}{30}% pseudocharacters-and-the-length-bound
\chapters@set{20261125T011000}{31}% galois-representations-black-box
\chapters@set{20261126T030300}{32}% the-data
\chapters@set{20261126T020210}{33}% bessel-poincare-series-and-their-gram-operators
\chapters@set{20261126T030400}{34}% the-seesaw-identity
\chapters@set{20261126T030500}{35}% divisibility-of-the-composite-lift
\chapters@set{20261126T030600}{36}% theorem-a-short-forms
\chapters@set{20261126T030700}{37}% the-determinant-of-a-chain
\chapters@set{20261126T030800}{38}% the-cost-of-one-lift
\chapters@set{20261126T030900}{39}% denominators-of-projections-to-lifts
\chapters@set{20261126T031000}{40}% twists-with-few-congruences
\chapters@set{20261126T031100}{41}% theorem-c-and-the-main-inequality
\providecommand\chapnumber[1]{%
  \@ifundefined{chapters@#1}%
    {\PackageError{chapters}{Sheet #1 is not listed in the index}%
      {List it in 20261125T012000--grh-workbook__guide.org\MessageBreak
       and run scripts/make-chapters.py.}}%
    {\edef\thechap{\@nameuse{chapters@#1}}}}
\makeatother
\chapnumber{<sheet ID>}
\makeatletter
\providecommand\chapters@set[2]{\@namedef{chapters@#1}{#2}}
\chapters@set{20261125T001000}{1}% zeros-to-ramp-sums
\chapters@set{20261125T001100}{2}% inert-primes-and-estimate-b
\chapters@set{20261125T005000}{3}% galois-fourier-and-averaging
\chapters@set{20261125T005500}{4}% columns-from-a-sphere
\chapters@set{20261125T000000}{5}% slope-method
\chapters@set{20261126T030100}{6}% a-twist-for-each-prime
\chapters@set{20261125T003000}{7}% the-chain-and-the-square-root
\chapters@set{20261125T010000}{8}% rank-lemma-and-zero-estimate
\chapters@set{20261125T010500}{9}% poincare-series-and-near-orthogonality
\chapters@set{20261124T223000}{10}% elliptic-curves-for-the-proof
\chapters@set{20261125T009500}{11}% fock-model-and-theta-intertwiners
\chapters@set{20261124T223100}{12}% complex-multiplication-and-inert-reduction
\chapters@set{20261125T008000}{13}% one-cm-determinant
\chapters@set{20261124T214500}{14}% supersingular-theta-divisibility
\chapters@set{20261124T214600}{15}% the-isotropic-plane
\chapters@set{20261125T013700}{16}% pairings-on-a-formal-group-of-height-two
\chapters@set{20261125T013500}{17}% hecke-nondegeneracy-on-the-tree
\chapters@set{20261126T030200}{18}% the-bound-is-attained
\chapters@set{20261125T009000}{19}% kronecker-limit-and-chowla-selberg
\chapters@set{20261125T002000}{20}% definite-theta-lifts
\chapters@set{20261125T015500}{21}% siegel-weil-in-the-smallest-case
\chapters@set{20261125T002100}{22}% definite-theta-lifts-and-norms
\chapters@set{20261125T016500}{23}% hecke-operators-and-theta-lifts
\chapters@set{20261125T016700}{24}% divisibility-along-the-chain
\chapters@set{20261125T004000}{25}% denominators-and-sha
\chapters@set{20261125T004100}{26}% tate-twists-and-sha
\chapters@set{20261126T020430}{27}% the-minus-part-by-gauss-sums
\chapters@set{20261125T014100}{28}% congruence-modules
\chapters@set{20261125T014500}{29}% eisenstein-congruences-and-selmer-classes
\chapters@set{20261125T014700}{30}% pseudocharacters-and-the-length-bound
\chapters@set{20261125T011000}{31}% galois-representations-black-box
\chapters@set{20261126T030300}{32}% the-data
\chapters@set{20261126T020210}{33}% bessel-poincare-series-and-their-gram-operators
\chapters@set{20261126T030400}{34}% the-seesaw-identity
\chapters@set{20261126T030500}{35}% divisibility-of-the-composite-lift
\chapters@set{20261126T030600}{36}% theorem-a-short-forms
\chapters@set{20261126T030700}{37}% the-determinant-of-a-chain
\chapters@set{20261126T030800}{38}% the-cost-of-one-lift
\chapters@set{20261126T030900}{39}% denominators-of-projections-to-lifts
\chapters@set{20261126T031000}{40}% twists-with-few-congruences
\chapters@set{20261126T031100}{41}% theorem-c-and-the-main-inequality
\providecommand\chapnumber[1]{%
  \@ifundefined{chapters@#1}%
    {\PackageError{chapters}{Sheet #1 is not listed in the index}%
      {List it in 20261125T012000--grh-workbook__guide.org\MessageBreak
       and run scripts/make-chapters.py.}}%
    {\edef\thechap{\@nameuse{chapters@#1}}}}
\makeatother
\chapnumber{20261125T004000}
\renewcommand{\thesection}{\thechap.\arabic{section}}

\theoremstyle{definition}
\newtheorem{problem}{Problem}[section]
\theoremstyle{plain}
\newtheorem{theorem}{Theorem}[section]

\newcommand{\Z}{\mathbb{Z}}
\newcommand{\Q}{\mathbb{Q}}
\newcommand{\R}{\mathbb{R}}
\newcommand{\Oc}{\mathcal{O}}
\newcommand{\Cl}{\operatorname{Cl}}
\newcommand{\Gal}{\operatorname{Gal}}
\newcommand{\len}{\operatorname{len}}
\newcommand{\lcm}{\operatorname{lcm}}
\newcommand{\num}{\operatorname{num}}

\setlength{\emergencystretch}{2em}
\title{Denominators and the Stickelberger element}
\author{}
\date{November 25, 2026}

\begin{document}
\maketitle
\refstepcounter{chap}\label{chap}

\paragraph{Prerequisites.}
Ideal class groups \cite[Ch.~I, \S\S3, 6]{Ne} and the norm of ideals in an extension of number fields.  Cyclotomic fields, Bernoulli numbers and generalized Bernoulli numbers \cite[Ch.~4--6]{Wa}.  Euclidean lattices, their covolumes and the indices of sublattices, as in Chapter~\ref{sm-chap}, \emph{Size against divisibility}.  Chapter~\ref{tt-chap}, \emph{Odd Tate twists}, continues this one.

\section{Denominators of linear maps}\label{sec:denominators}

A linear map between \(\Q\)-vector spaces with integral structures need not respect them.  What does that cost when size is compared with divisibility, and how can it be reduced?  The simplest example is differentiation.  Let \(\mathrm{Int}\) be the set of polynomials \(f \in \Q[x]\) with \(f(\Z) \subseteq \Z\), and let \(\mathrm{Int}_n\) be its elements of degree at most \(n\).

\begin{problem}\label{pr:intval}
The binomial polynomials \(\binom{x}{k}\), \(0 \le k \le n\), form a \(\Z\)-basis of \(\mathrm{Int}_n\).
\end{problem}

\begin{problem}\label{pr:derivative}
Let \(D = d/dx\).  Then
\[
  D \binom{x}{n} = \sum_{k=1}^{n} \frac{(-1)^{k-1}}{k} \binom{x}{n-k} .
\]
Hence the least positive integer \(c\) with \(c \, D(\mathrm{Int}_n) \subseteq \mathrm{Int}\) is \(\lcm(1, 2, \dots, n)\).
\end{problem}

The logarithm of this denominator is \(\sum_{p^a \le n} \log p\), of the size of \(n\).  Alpöge's note records a ``need for division'' \cite[item~14]{Alp}.  In Alp\"oge's proof the denominator that matters is that of the orthogonal projection of integral forms onto a space of theta lifts, and a congruence module measures it (Chapter~\ref{cg-chap}, \emph{Congruence modules}).

An integral structure can be built one prime at a time.  For a lattice \(L \subset \Q^r\) and a prime \(p\), write \(L_p := L \otimes \Z_p \subset \Q_p^r\).

\begin{problem}\label{pr:local-global}
For every family \((\Lambda_p)_p\) of \(\Z_p\)-lattices \(\Lambda_p \subset \Q_p^r\) with \(\Lambda_p = \Z_p^r\) for all but finitely many \(p\), there is exactly one lattice \(L \subset \Q^r\) with \(L_p = \Lambda_p\) for every \(p\), namely \(L = \bigcap_p (\Lambda_p \cap \Q^r)\).  If \(\Lambda_2 = \tfrac12 \Z_2 \oplus \Z_2\), \(\Lambda_3 = \Z_3 \oplus \tfrac13 \Z_3\), and \(\Lambda_p = \Z_p^2\) otherwise, then \(L = \tfrac12\Z \oplus \tfrac13\Z\).
\end{problem}

For a finitely generated torsion \(\Z_p\)-module \(T\) let \(\len_{\Z_p} T\) be its length.

\begin{problem}\label{pr:length}
If \(T\) is a finite abelian group, then \(\log \lvert T \rvert = \sum_p \len_{\Z_p}(T \otimes \Z_p) \log p\).  If \(M' \subseteq M\) are lattices of the same rank, then \(\log [M : M'] = \sum_p \len_{\Z_p}(M_p/M'_p) \log p\).
\end{problem}

Let \(L\) and \(M\) be Euclidean lattices and \(T \colon L_\Q \to M_\Q\) an injective linear map.  If \(T(L) \not\subseteq M\), one may replace \(M\) by the lattice \(M + T(L)\), with the same norm.

\begin{problem}\label{pr:enlarge}
Suppose that \(L\) and \(M\) have the same rank.  Then \(\operatorname{covol}(M + T(L)) = \operatorname{covol}(M)/[M + T(L) : M]\), and
\[
  \operatorname{covol}(M) \le [M + T(L) : M]\, \lvert\det T\rvert\, \operatorname{covol}(L),
\]
where \(\det T\) is taken in orthonormal bases of \(L \otimes \R\) and \(M \otimes \R\).
\end{problem}

So a denominator costs \(\log\lvert(M + T(L))/M\rvert\), a sum of local terms by Problem~\ref{pr:length}.  If several maps \(T_1, \dots, T_k\) serve equally well, for instance because they have the same period, one may choose at each prime the lattice \((M + T_j(L))_p\) of smallest index over \(M_p\), and Problem~\ref{pr:local-global} assembles the choices into one lattice \(M^*\).  A map carrying \(L\) into \(M^*\) must then be found separately; for the Stickelberger element below it is a combination of the twists.

\begin{problem}\label{pr:gcd}
Let \(M_j := M + T_j(L)\), and let \(M^*\) be the lattice with \(M^*_p = (M_j)_p\) for a \(j = j(p)\) that minimizes \(\len_{\Z_p}\bigl((M_j)_p/M_p\bigr)\).  Then \(M^* \supseteq M\) and
\[
  \log [M^* : M] = \sum_p \min_j \len_{\Z_p}\bigl((M_j)_p/M_p\bigr) \log p .
\]
If \(M = \Z\) and \(M_j = D_j^{-1}\Z\) with positive integers \(D_j\), then \([M^* : M] = \gcd(D_1, \dots, D_k)\).  For \(D_j = 6, 10, 15\) each map alone costs at least \(\log 6\), and together they cost nothing.
\end{problem}

Alpöge's note glues its integral structures in this way, ``across Hecke characters for a `gcd' '' \cite[item~15]{Alp}.  What no twist removes is, at each prime, a finite module that it identifies with a group of Galois cohomology classes \cite[item~17]{Alp}; in Alp\"oge's proof these modules are bounded through Selmer groups of twists, and what remains is a group of odd classes, bounded at every prime by one integer (Chapter~\ref{oc-chap}, \emph{Odd classes by Gauss sums}).  The odd classes are the minus part, which is what the Stickelberger element below sees.  The simplest instance, a rational element whose twists are integral and whose remainder is measured by a class group, is the Stickelberger element of \(\Q(\zeta_p)\).

\paragraph{In the Lean formalization.}
In the rewritten Lean proof \cite{GL} the denominator of the projection is \texttt{ClearsProjection}, in \texttt{New/Determinant/Twist/FitTwists.lean} of the readable edition.  For a twist, an isomorphism \(\iota \colon \mathbb{C} \simeq \bar\Q_\ell\), the space \(Z\) of lifts from one definite space to the next and a real \(\beta\), it asks for an operator \(\Lambda\) on \(Z\) with \(\lvert\iota(\det\Lambda)\rvert \ge e^{-r\beta\dim Z}\), where \(r\) is a rank of the proof, such that \(\Lambda(\mathrm{pr}_Z\,y)\) is integral at \(\iota\) for every form \(y\) integral at \(\iota\) and orthogonal to certain false lifts.  The sum over \(\ell\) of the numbers \(\beta\) must stay within a budget (\texttt{IsDenominatorBound}).  The odd classes are bounded by \texttt{New.exists\_length\_shaOdd\_le}, in \texttt{Lib/Cyclotomic/OddSha.lean}.

\section{The Stickelberger element}\label{sec:stickelberger}

Let \(p\) be an odd prime, \(F := \Q(\zeta_p)\), \(G := \Gal(F/\Q)\), and let \(\sigma_a \in G\) be the automorphism with \(\sigma_a(\zeta_p) = \zeta_p^a\).  Let \(A\) be the \(p\)-primary part of \(\Cl(F)\), a \(\Z_p[G]\)-module.  Let \(\omega \colon (\Z/p\Z)^\times \to \Z_p^\times\) be the Teichmüller character, with \(\omega(a) \equiv a \pmod p\), and view it as a character of \(G\).  For \(i \in \Z/(p-1)\Z\) let \(A^{(i)} := \{ z \in A : \sigma_a z = \omega(a)^i z \text{ for all } a \}\), and
\[
  \varepsilon_i := \frac{1}{p-1} \sum_{a=1}^{p-1} \omega(a)^i \sigma_a^{-1} \in \Z_p[G] .
\]

\begin{problem}\label{pr:eigen}
The \(\varepsilon_i\) are orthogonal idempotents with sum \(1\), \(\sigma_a \varepsilon_i = \omega(a)^i \varepsilon_i\), and \(A = \bigoplus_i \varepsilon_i A = \bigoplus_i A^{(i)}\).  Complex conjugation \(\sigma_{-1}\) acts on \(A^{(i)}\) by \((-1)^i\).
\end{problem}

Write \(A^+\) and \(A^-\) for the sums of the \(A^{(i)}\) with \(i\) even and with \(i\) odd.

\begin{problem}\label{pr:subfield}
If \(K \subseteq F\), then extension of ideals identifies the \(p\)-primary part of \(\Cl(K)\) with \(A^{\Gal(F/K)}\).  Hence \(A^{(0)} = 0\), \(A^+\) is the \(p\)-primary part of the class group of \(F^+ := \Q(\zeta_p + \zeta_p^{-1})\), and if \(p \equiv 1 \pmod 4\), then \(A^{((p-1)/2)}\) is the \(p\)-primary part of \(\Cl(\Q(\sqrt{p}))\).
\end{problem}

The \emph{Stickelberger element} is
\[
  \theta := \frac1p \sum_{a=1}^{p-1} a \, \sigma_a^{-1} \in \Q[G] ,
\]
with the denominator \(p\).  We use Stickelberger's theorem as a black box \cite[Thm.~6.10]{Wa}.

\begin{theorem}[Stickelberger]\label{thm:stickelberger}
If \(\beta \in \Z[G]\) and \(\beta\theta \in \Z[G]\), then \(\beta\theta\) annihilates \(\Cl(F)\).
\end{theorem}

The integral multiples of \(\theta\) are its twists.

\begin{problem}\label{pr:c-twist}
For \(c\) prime to \(p\), \((c - \sigma_c)\theta = \sum_{b=1}^{p-1} \lfloor cb/p \rfloor \, \sigma_b^{-1}\).  In particular \((c - \sigma_c)\theta \in \Z[G]\).
\end{problem}

For a nontrivial Dirichlet character \(\chi\) modulo \(p\) let \(B_{1,\chi} := \frac1p \sum_{a=1}^{p-1} \chi(a)\, a\).

\begin{problem}\label{pr:eigen-theta}
\(\varepsilon_i \theta = B_{1,\omega^{-i}} \varepsilon_i\) for \(i \not\equiv 0\).  If \(i \not\equiv 0\) is even, then \(B_{1,\omega^{-i}} = 0\).
\end{problem}

So \(\theta\) says nothing about \(A^+\).  On \(A^{(i)}\) the twist \((c - \sigma_c)\theta\) acts as \((c - \omega(c)^i)B_{1,\omega^{-i}}\), and the gluing of Problem~\ref{pr:gcd} asks for the ideal generated by the factors \(c - \omega(c)^i\).

\begin{problem}\label{pr:gcd-c}
The ideal of \(\Z_p\) generated by the numbers \(c - \omega(c)^i\), \(c\) prime to \(p\), is \(\Z_p\) if \(i \not\equiv 1\) and \(p\Z_p\) if \(i \equiv 1\).  Moreover \(v_p(B_{1,\omega^{-1}}) = -1\).
\end{problem}

Thus on each \(A^{(i)}\) with \(i \not\equiv 0, 1\) the twists together act as \(B_{1,\omega^{-i}} \in \Z_p\): the gluing recovers the eigenvalue of \(\theta\) itself.  On \(A^{(1)}\) every factor \(c - \omega(c)\) is divisible by \(p\) and cancels the denominator of \(B_{1,\omega^{-1}}\), so the twists generate \(\Z_p\) and not \(p^{-1}\Z_p\): the gluing does not recover the eigenvalue of \(\theta\).  This \(p\) is the order of \(\mu_p\), on which \(G\) acts by \(\omega\).  (The group \(A^{(1)}\) vanishes for another reason \cite[Prop.~6.16]{Wa}.)  By Kummer's congruences \cite[Ch.~5]{Wa}, if \(i\) is odd and \(3 \le i \le p-2\), then
\begin{equation}\label{eq:kummer}
  B_{1,\omega^{-i}} \equiv \frac{B_{p-i}}{p-i} \pmod p .
\end{equation}

\begin{problem}\label{pr:p7}
For \(p = 7\) the character \(\omega^{-3}\) is the Legendre symbol, \(B_{1,\omega^{-3}} = -1\), and \(B_4/4 = -1/120 \equiv -1 \pmod 7\).
\end{problem}

\begin{problem}\label{pr:herbrand}
Let \(i\) be odd with \(3 \le i \le p-2\).  If \(p \nmid B_{p-i}\), then \(A^{(i)} = 0\).
\end{problem}

This is Herbrand's theorem \cite[Thm.~6.17]{Wa}.  Its converse is due to Ribet \cite{Ri}, and the exact order to Mazur and Wiles \cite{MW}; see \cite[Ch.~15]{Wa}.  We use them as a black box.

\begin{theorem}[Herbrand, Ribet, Mazur--Wiles]\label{thm:HR}
Let \(i\) be odd with \(3 \le i \le p-2\).  Then \(A^{(i)} \ne 0\) if and only if \(p \mid B_{p-i}\), and \(\lvert A^{(i)} \rvert = \lvert B_{1,\omega^{-i}} \rvert_p^{-1}\).
\end{theorem}

\begin{problem}\label{pr:irregular}
For \(p < 100\), the pairs \((p, k)\) with \(k\) even, \(2 \le k \le p-3\) and \(p \mid B_k\) are \((37, 32)\), \((59, 44)\) and \((67, 58)\).  Hence for \(p < 100\) the group \(A^-\) is nonzero only for \(p = 37, 59, 67\), and then exactly one odd eigenspace is nonzero, \(A^{(5)}\), \(A^{(15)}\) and \(A^{(9)}\) respectively.
\end{problem}

The plus part \(A^+\) is not seen by \(\theta\) (Problem~\ref{pr:eigen-theta}).  Its order is the \(p\)-part of an index of cyclotomic units in the units of \(F^+\) \cite[Thm.~8.2]{Wa}.  Vandiver's conjecture, that \(A^+ = 0\) for every \(p\), has not been proved; it has been verified for \(p < 2^{31}\) \cite{HHO}.  The next chapter meets the same division between minus and plus parts in Galois cohomology; the proof needs only the minus part.

\section*{Hints}

\subsection*{\S\ref{sec:denominators}}
\begin{description}
\item[\ref{pr:intval}] Use the difference operator \(\Delta f(x) := f(x+1) - f(x)\) and \(\Delta \binom{x}{k} = \binom{x}{k-1}\).
\item[\ref{pr:derivative}] Formally \(D = \log(1 + \Delta)\).  For the sharpness, use the uniqueness in Problem~\ref{pr:intval}.
\item[\ref{pr:local-global}] Choose an integer \(N\) with \(N \Z_p^r \subseteq \Lambda_p \subseteq N^{-1}\Z_p^r\) for all \(p\), and work in \(N^{-1}\Z^r / N\Z^r\).
\item[\ref{pr:length}] The \(p\)-primary part of \(T\) is \(T \otimes \Z_p\).
\item[\ref{pr:enlarge}] \(M\) and \(T(L)\) are sublattices of finite index of \(M + T(L)\).
\item[\ref{pr:gcd}] Apply Problem~\ref{pr:length} to \(M^*/M\).
\end{description}

\subsection*{\S\ref{sec:stickelberger}}
\begin{description}
\item[\ref{pr:eigen}] Orthogonality of the characters \(\omega^i\) of the cyclic group \(G\), whose order is prime to \(p\).
\item[\ref{pr:subfield}] Compose extension of ideal classes with the norm, in both orders; \([F:K]\) is prime to \(p\).  For the last claim, find the subfield fixed by the squares in \(G\).
\item[\ref{pr:eigen-theta}] Use \(\sigma_a^{-1}\varepsilon_i = \omega(a)^{-i}\varepsilon_i\).  For even \(\chi\), pair \(a\) with \(p - a\).
\item[\ref{pr:gcd-c}] Reduce \(c - \omega(c)^i\) modulo \(p\).  For \(B_{1,\omega^{-1}}\), reduce \(p B_{1,\omega^{-1}}\) modulo \(p\).
\item[\ref{pr:herbrand}] Apply Theorem~\ref{thm:stickelberger} to \(\beta = c - \sigma_c\), project with \(\varepsilon_i\), choose \(c\) by Problem~\ref{pr:gcd-c}, and use~\eqref{eq:kummer}.
\item[\ref{pr:irregular}] The first part is a finite computation.  For the second use Theorem~\ref{thm:HR} with \(k = p - i\), and \(A^{(1)} = 0\) \cite[Prop.~6.16]{Wa}.
\end{description}

\begin{thebibliography}{HHO}
\bibitem[Alp]{Alp} L.~Alp\"oge, untitled note announcing a proof of the Riemann hypothesis, 20 November 2026, \url{https://alpo.ge/remark.pdf}.
\bibitem[GL]{GL} L.~Sernau, \emph{GRH for Dirichlet \(L\)-functions, in Lean~4}, a shorter rewriting of the Lean proof, text by Claude, \url{https://github.com/lsernau/grh-lean}, commit \texttt{72fead1}, readable edition \texttt{grh/}.
\bibitem[HHO]{HHO} W.~Hart, D.~Harvey and W.~Ong, \emph{Irregular primes to two billion}, Math.\ Comp.\ \textbf{86} (2017), 3031--3049.
\bibitem[MW]{MW} B.~Mazur and A.~Wiles, \emph{Class fields of abelian extensions of \(\mathbf{Q}\)}, Invent.\ Math.\ \textbf{76} (1984), 179--330.
\bibitem[Ne]{Ne} J.~Neukirch, \emph{Algebraic Number Theory}, Grundlehren Math.\ Wiss.\ 322, Springer, 1999.
\bibitem[Ri]{Ri} K.~A.~Ribet, \emph{A modular construction of unramified \(p\)-extensions of \(\mathbf{Q}(\mu_p)\)}, Invent.\ Math.\ \textbf{34} (1976), 151--162.
\bibitem[Wa]{Wa} L.~C.~Washington, \emph{Introduction to Cyclotomic Fields}, 2nd ed., Grad.\ Texts in Math.\ 83, Springer, 1997.
\end{thebibliography}

\end{document}
